% !TEX root = ../main.tex
What evidence distinguishes collusion from a mistake, an unexecuted promise,
or an inability to complete a task? We study five model families in a dispute
game where scripted parties offer jurors conditional rewards for accepting a
false claim or rejecting a supported claim. The original frame crosses five
launch-seed labels and six conditions on two fixed finance and education
sources. Among 300 assigned cases, the frozen record contains 122 source-correct
settlements, two source-adverse settlements, 26 aborts and 150 unresolved
outcomes. Both wrong settlements are GPT OSS education cases with executed
extra payments of 10 synthetic ETH to each of two jurors. Returned text
explicitly considers reward eligibility, but knowledge of falsity is not
established. One paid wrong settlement occurs inside a failed run, demonstrating
why case evidence must survive session-level failure. Selected recovery cases
contrast Inkling's explicit refusal and correct execution with DeepSeek's
correct arithmetic without a commitment. These observations motivate separate
labels for coordination, implementing action, source knowledge and capability.
The separately registered exploratory recovery is not pooled with confirmation.
Six conditions share one paid-offer amount; they are not a stakes sweep.
We provide a five-model mechanical outcome profile, not a validated
collusion-resistance ranking: independent semantic review remains open, so
resistance estimates and confidence intervals are unavailable. The planned
detector comparison will test communication, action and combined evidence
against fixed independent labels.
